\section{Palindrome Partitioning}
\label{sec:rec-palindrome-partitioning}

\problemstatement{Given a string $s$, partition it into one or more
contiguous substrings such that \textbf{every} substring in the
partition is itself a palindrome. Return every such partitioning.}

\begin{patternbox}
\emph{``Partition a sequence into pieces, each satisfying some
property''} introduces a new choice shape: instead of choosing
\emph{which elements} to include (as in subsets/combinations) or
\emph{which direction} to move (as in grid backtracking), the choice here
is \textbf{how far the next piece should extend} -- a loop over every
possible end position for the current piece, starting from wherever the
previous piece left off.
\end{patternbox}

\subsection*{Understanding the problem carefully}

Starting from some position $\textit{start}$ in the string, the next
piece of the partition can end anywhere from $\textit{start}$ itself up
to the end of the string -- but only positions where the resulting
substring $s[\textit{start}..\textit{end}]$ is a palindrome are valid
choices. For each valid choice of where this piece ends, the rest of the
string (from $\textit{end}+1$ onward) must be recursively partitioned
the same way; reaching the very end of the string with no characters
left to partition marks one complete, valid partitioning.

\begin{ideabox}[Loop Over Every Valid Palindromic Prefix from the Current Start]
$\textsc{Solve}(\textit{start}, \textit{current})$: if
$\textit{start} = |s|$: record a copy of $\textit{current}$; return.
Otherwise: for $\textit{end}$ from $\textit{start}$ to $|s|-1$: if
$s[\textit{start}..\textit{end}]$ is a palindrome: append this
substring to $\textit{current}$; recurse
$\textsc{Solve}(\textit{end}+1, \textit{current})$; remove the
substring (undo).
\end{ideabox}

\subsection*{Why looping over end positions correctly enumerates every partitioning}

Every valid partitioning of $s[\textit{start}..]$ is uniquely determined
by where its \emph{first} piece ends -- once that is fixed (to some
valid palindromic substring), the rest of the partitioning is exactly a
partitioning of whatever remains, which is the same sub-problem at a
smaller starting position. By trying \emph{every} valid place the first
piece could end, and recursively solving the remainder for each choice,
we generate every possible partitioning exactly once, with no overlap
between different choices of where the first piece ends (since two
different end positions for the first piece necessarily produce
different first pieces, hence different overall partitionings).

\subsection*{Algorithm}

\begin{enumerate}
  \item $\textsc{Solve}(\textit{start}, \textit{current})$:
  \begin{enumerate}
    \item If $\textit{start} = |s|$: record $\textit{current}$; return.
    \item For $\textit{end}$ from $\textit{start}$ to $|s|-1$:
    \begin{itemize}
      \item If $s[\textit{start}..\textit{end}]$ is a palindrome:
            append it to $\textit{current}$; recurse
            $\textsc{Solve}(\textit{end}+1, \textit{current})$; remove
            it (undo).
    \end{itemize}
  \end{enumerate}
  \item Call $\textsc{Solve}(0, [\,])$.
\end{enumerate}

\subsection*{Dry run}

\dryrun{Take $s = \texttt{"aab"}$.}

\begin{itemize}
  \item $\textsc{Solve}(0,[])$: $\textit{end}=0$:
        $s[0..0]=\texttt{"a"}$, palindrome. Append. Recurse
        $\textsc{Solve}(1,[\texttt{"a"}])$.
  \begin{itemize}
    \item $\textit{end}=1$: $s[1..1]=\texttt{"a"}$, palindrome.
          Append. Recurse $\textsc{Solve}(2,[\texttt{"a"},\texttt{"a"}])$.
    \begin{itemize}
      \item $\textit{end}=2$: $s[2..2]=\texttt{"b"}$, palindrome.
            Append. Recurse $\textsc{Solve}(3,[\texttt{"a"},\texttt{"a"},\texttt{"b"}])$:
            $\textit{start}=3=|s|$, record
            $[\texttt{"a"},\texttt{"a"},\texttt{"b"}]$.
    \end{itemize}
    \item Undo. $\textit{end}=2$ at the $[\texttt{"a"}]$ level:
          $s[1..2]=\texttt{"ab"}$, not a palindrome. Skip.
  \end{itemize}
  \item Undo back to $[]$. $\textit{end}=1$ at the top level:
        $s[0..1]=\texttt{"aa"}$, palindrome. Append. Recurse
        $\textsc{Solve}(2,[\texttt{"aa"}])$.
  \begin{itemize}
    \item $\textit{end}=2$: $s[2..2]=\texttt{"b"}$, palindrome.
          Append. Recurse $\textsc{Solve}(3,[\texttt{"aa"},\texttt{"b"}])$:
          record $[\texttt{"aa"},\texttt{"b"}]$.
  \end{itemize}
  \item $\textit{end}=2$ at the top level: $s[0..2]=\texttt{"aab"}$,
        not a palindrome. Skip.
\end{itemize}

Recorded partitionings: $[\texttt{"a"},\texttt{"a"},\texttt{"b"}]$ and
$[\texttt{"aa"},\texttt{"b"}]$ -- the only two ways to partition
\texttt{"aab"} into all-palindrome pieces.

\subsection*{C++ implementation}

\begin{lstlisting}
#include <bits/stdc++.h>
using namespace std;

bool isPalindrome(string& s, int lo, int hi) {
    while (lo < hi) {
        if (s[lo] != s[hi]) return false;
        lo++; hi--;
    }
    return true;
}

void solve(int start, string& s, vector<string>& current,
           vector<vector<string>>& result) {
    if (start == (int)s.size()) {
        result.push_back(current);
        return;
    }
    for (int end = start; end < (int)s.size(); end++) {
        if (isPalindrome(s, start, end)) {
            current.push_back(s.substr(start, end - start + 1));
            solve(end + 1, s, current, result);
            current.pop_back();
        }
    }
}

vector<vector<string>> palindromePartitioning(string s) {
    vector<vector<string>> result;
    vector<string> current;
    solve(0, s, current, result);
    return result;
}

int main() {
    for (auto& partition : palindromePartitioning("aab")) {
        cout << "[ ";
        for (string& piece : partition) cout << piece << " ";
        cout << "] ";
    }
    cout << endl;
    return 0;
}
\end{lstlisting}

\begin{complexitybox}
\textbf{Time Complexity.} Up to $O(2^n)$ partitionings can exist in the
worst case (e.g.\ a string of all-identical characters, where every
position could start a new piece), and checking each candidate substring
for the palindrome property costs up to $O(n)$, giving an overall worst
case of
\[
  O(n \cdot 2^n).
\]
\textbf{Space Complexity.} $O(n)$ for the recursion depth.
\end{complexitybox}

\begin{takeawaybox}
``Partition a sequence into pieces satisfying a property'' is a third
fundamental backtracking choice shape, alongside ``include or exclude
each element'' and ``move in one of several directions'': loop over
every possible \emph{extent} of the next piece from the current
position, recurse on the remainder for each valid extent. This same
``loop over partition points'' shape reappears, generalized with memoization,
in the Matrix Chain Multiplication / Partition DP family covered later in
the dynamic programming chapters.
\end{takeawaybox}
